\section{Atividades}

\subsection{Atividade 1 -- Análise de Relatórios Públicos}

\paragraph{}
Inicialmente foram analisados relatórios públicos de testes de invasão
disponibilizados por profissionais e organizações da área de segurança da
informação.

\paragraph{}
Durante a análise foram observadas as estruturas utilizadas nos documentos,
incluindo seções como escopo, metodologia, sumário executivo, descrição das
vulnerabilidades encontradas, evidências e recomendações de mitigação.

\paragraph{}
Também foi possível identificar diferentes formas de apresentação das
informações e compreender como os relatórios procuram comunicar os resultados
tanto para equipes técnicas quanto para gestores.

\subsection{Atividade 2 -- Coleta de Evidências Estáticas}

\paragraph{}
Nesta atividade foram utilizadas ferramentas voltadas à coleta de evidências
durante a execução de testes de segurança.

\paragraph{}
Foi explorado o comando \texttt{script}, responsável por registrar comandos e
saídas de terminais em arquivos de texto, permitindo manter um histórico
completo das ações realizadas durante um processo de avaliação.

\paragraph{}
Também foi utilizada a ferramenta \texttt{flameshot}, empregada para captura e
anotação de imagens, possibilitando destacar informações relevantes por meio de
setas, marcações, textos e outros recursos visuais.

\paragraph{}
A atividade demonstrou a importância da documentação contínua durante a
execução dos testes para facilitar a elaboração posterior dos relatórios.

\subsection{Atividade 3 -- Coleta de Evidências Dinâmicas}

\paragraph{}
Posteriormente foram avaliadas ferramentas destinadas à geração de evidências
dinâmicas, permitindo registrar comportamentos observados ao longo do tempo.

\paragraph{}
Foi utilizada a ferramenta \texttt{peek} para gravação de GIFs contendo
demonstrações rápidas de funcionalidades e comportamentos identificados durante
os testes.

\paragraph{}
Além disso, foi explorado o uso da ferramenta \texttt{RecordMyDesktop} para
captura de vídeos completos das atividades realizadas no ambiente.

\paragraph{}
Esses recursos possibilitam demonstrar ataques, explorações e validações de
forma mais clara quando comparados exclusivamente a imagens estáticas.

\subsection{Atividade 4 -- Coleta de Evidências de Ataque}

\paragraph{}
Nesta etapa foi realizada novamente a atividade prática relacionada a ataques
de força bruta, desta vez com foco específico na documentação e preservação das
evidências geradas durante a execução dos testes.

\paragraph{}
Foram registrados comandos executados, resultados obtidos, credenciais
identificadas e demais informações relevantes utilizando as ferramentas
apresentadas anteriormente.

\paragraph{}
Também foram coletadas capturas de tela e registros complementares para apoiar
a elaboração do relatório técnico final.

\paragraph{}
A atividade permitiu compreender a importância de registrar adequadamente todas
as etapas de um processo de avaliação de segurança.

\subsection{Atividade 5 -- Escrita de Relatório}

\paragraph{}
Por fim, foi realizada a elaboração de um relatório técnico baseado nas
evidências coletadas durante as atividades anteriores.

\paragraph{}
O documento produzido descreveu as etapas executadas, os testes realizados, os
resultados obtidos e as vulnerabilidades identificadas durante a atividade
prática.

\paragraph{}
Também foram incluídas evidências, descrições técnicas e recomendações de
mitigação, seguindo as boas práticas observadas durante a análise dos
relatórios públicos.

\paragraph{}
Essa atividade consolidou os conhecimentos relacionados à documentação técnica
e à comunicação adequada dos resultados de avaliações de segurança.
